\section{Conclusion}
\label{sec:conclusion}

The goal of our work is to assay fairness in machine learning within the context
of causal reasoning. This perspective naturally addresses shortcomings of
earlier statistical approaches. Causal fairness criteria are suitable whenever
we are willing to make assumptions about the (causal) generating process
governing the data.  Whilst not always feasible, the causal approach naturally
creates an incentive to scrutinize the data more closely and work out plausible
assumptions to be discussed alongside any conclusions regarding fairness.

Key concepts of our conceptual framework are \emph{resolving variables} and
\emph{proxy variables} that play a dual role in defining causal discrimination
criteria. We develop a practical procedure to remove proxy discrimination given
the structural equation model and analyze a similar approach for unresolved
discrimination. In the case of proxy discrimination for linear structural
equations, the procedure has an intuitive form that is similar to heuristics
already used in the regression literature. Our framework is limited by the
assumption that we can construct a valid causal graph. The removal of proxy
discrimination moreover depends on the functional form of the causal
dependencies. We have focused on the conceptual and theoretical analysis, and
experimental validations are beyond the scope of the present work.

The causal perspective suggests a number of interesting new directions at the
technical, empirical, and conceptual level. We hope that the framework and
language put forward in our work will be a stepping stone for future
investigations.
